% This must be in the first 5 lines to tell arXiv to use pdfLaTeX, which is strongly recommended.
\pdfoutput=1
% In particular, the hyperref package requires pdfLaTeX in order to break URLs across lines.

\documentclass[11pt]{article}

% Remove the "review" option to generate the final version.
\usepackage[review]{acl}

% Standard package includes
\usepackage{times}
\usepackage{latexsym}

\usepackage{booktabs}
\usepackage{multirow} 
\usepackage{makecell}

% For proper rendering and hyphenation of words containing Latin characters (including in bib files)
\usepackage[T1]{fontenc}
% For Vietnamese characters
% \usepackage[T5]{fontenc}
% See https://www.latex-project.org/help/documentation/encguide.pdf for other character sets

% This assumes your files are encoded as UTF8
\usepackage[utf8]{inputenc}

% This is not strictly necessary, and may be commented out.
% However, it will improve the layout of the manuscript,
% and will typically save some space.
\usepackage{microtype}

% This is also not strictly necessary, and may be commented out.
% However, it will improve the aesthetics of text in
% the typewriter font.
\usepackage{inconsolata}

%Including images in your LaTeX document requires adding
%additional package(s)
\usepackage{graphicx}

% If the title and author information does not fit in the area allocated, uncomment the following
%
%\setlength\titlebox{<dim>}
%
% and set <dim> to something 5cm or larger.

\title{Team Cantharellus at SemEval-2025 Task 3: Multilingual Detection of Hallucinations in LLM Outputs}

% Author information can be set in various styles:
% For several authors from the same institution:
 \author{Xinyan Mo \and Nikolay Vorontsov \and Tiankai Zang \\
         \{xinyan.mo, nikolay.vorontsov, tiankai.zang\}@helsinki.fi}
% if the names do not fit well on one line use
%         Author 1 \\ {\bf Author 2} \\ ... \\ {\bf Author n} \\
% For authors from different institutions:
% \author{Author 1 \\ Address line \\  ... \\ Address line
%         \And  ... \And
%         Author n \\ Address line \\ ... \\ Address line}
% To start a seperate ``row'' of authors use \AND, as in
% \author{Author 1 \\ Address line \\  ... \\ Address line
%         \AND
%         Author 2 \\ Address line \\ ... \\ Address line \And
%         Author 3 \\ Address line \\ ... \\ Address line}

\author{First Author \\
  Affiliation / Address line 1 \\
  Affiliation / Address line 2 \\
  Affiliation / Address line 3 \\
  \texttt{email@domain} \\\And
  Second Author \\
  Affiliation / Address line 1 \\
  Affiliation / Address line 2 \\
  Affiliation / Address line 3 \\
  \texttt{email@domain} \\}

\begin{document}
\maketitle
\begin{abstract}
This paper describes our submission to SemEval-2025 Task-3: Mu-shroom, the Multilingual Shared-task on Hallucinations and Related Observable Overgeneration Mistakes, which mainly aims at detecting spans of LLM-generated text corresponding to hallucinations in multilingual and multimodel context. We propose two distinctive approaches: 1) fine-tuning pretrained language models and 2) prompting generative large language models. The results show that {to be filled}. We participated in 13 out of 14 available languages and achieved an average ranking of 10th out of 41 participating teams, with our highest ranking securing a top 5 position.


\end{abstract}


\section{Introduction}
Recent years have witnessed the rapid development of large language models (LLMs) and their applications in various fields of natural language processing \cite{wei2022chain, zhao2023survey}. However, content generated by LLMs occasionally contains inaccurate or fictitious information \cite{perkovic2024hallucinations}. The phenomenon in which natural language generation models often generate text that is nonsensical, or unfaithful to the provided source input is commonly referred to as “hallucinations” \cite{ji2023survey}. It is therefore a vital task to detect and identify hallucinated content so as to improve the reliability and trustworthiness of LLM-generated content.

SemEval 2025 Task 3 (Mu-SHROOM, the Multilingual Shared-task on Hallucinations and Related Observable Overgeneration Mistakes) (some citation) proposes the task of detecting hallucination in content generated with LLMs. Different from the previous iteration, SemEval-2024 Task 6 \cite{mickus2024semeval}, in which the participants were asked to make binary decisions of whether a given context contains hallucination, the current task requires the participants to predict where the hallucinations occur. Specifically, the current task requires the participants to predict the spans of the hallucinated content within LLM outputs in 14 different languages.
This paper describes our participation in the task in detail. We participated in 13 out of 14 languages, and introduced two distinctive approaches of hallucination span prediction.

\textbf{1. Fine-tuning pretrained language models:}  As is shown in the results of the previous iteration of this task \cite{mickus2024semeval}, it is effective to fine tune pretrained language model for hallucination detection. We therefore further extended this approach from performing binary classification tasks to predict the spans of hallucinations. We explored fine-tuning a series of Transformer-based pretrained language models \cite{vaswani2017attention}, including text-to-text transformer models \cite{raffel2020exploring} and Bert-based models \cite{devlin2018bert}, with training data created by ourselves.

\textbf{2. Prompting generative LLMs:} Current state-of-the-art LLM such as GPT-4 \cite{achiam2023gpt} has exhibited impressive abilities in a wide range of tasks. As the third approach, we proposed prompting such models to generate predictions directly. To that end, we explored various prompting techniques on a GPT-4o \cite{hurst2024gpt}, the current state-of-the-art model for natural language generation.
We have released our code and other relevant material on GitHub. 
[footnote and link here]


\section{Background}
Hallucination detection has been an extensively researched topic in recent years. One promising solution for hallucination detection is to utilize the self-evaluation ability of LLMs to judge the factual correctness of a given statement, leveraging the fact that LLMs have possessed a rich knowledge base \cite{li2024dawn, zhang2024self}. Many existing studies focus on the binary classification of hallucination. For instance, SelfCheckGPT \cite{manakul2023selfcheckgpt} is built on the idea that when an LLM is familiar with a particular concept, the responses it generates are likely to be consistent and contain similar facts. In the HaluEval 2.0 benchmark \cite{li2024dawn}, the hallucination detection approach is built by first extracting factual statements from LLM responses, then determining the trustfulness of these statements with respect to world knowledge by taking advantage of the vast knowledge base of LLMs. Similarly, MIND \cite{su2024unsupervised} introduces a similar approach that leverages the internal states of LLMs for real-time hallucination detection without requiring manual annotations. GraphEval \cite{sansford2024grapheval} proposes a method of detecting inconsistencies with respect to provided knowledge using a knowledge-graph approach.

Prompt engineering is a crucial technique for extending the capabilities of LLMs \cite{sahoo2024systematic}. A commonly used strategy is few-shot prompting, which refers to the technique of constructing prompts using a small set of demonstrative input-output examples \cite{lazaridou2022internet, ma2023fairness}. Another method that has been proved effective is chain-of-thought (CoT) prompting, which is achieved by guiding the LLM to think step by step to perform complex reasoning tasks \cite{wei2022chain, zhang2022automatic, chen2023you}.

Named entity recognition (NER) is the task of identifying named entities such as person, location and organization in text \cite{yadav2019survey}. Deep learning approaches have increasingly been adopted for NER and have exhibited impressive abilities across various domains and languages \cite{li2020survey, song2021deep, liu2022chinese}. Recent research also explores the possibility of integrating prompting techniques into NER tasks by utilizing LLMs \cite{shen2023promptner, hu2024improving}.

\section{System Overview}
\subsection{The Prompt Engineering Approach}
[add content]

\subsection{The Fine-Tuning Approach}
The goal of this task is to detect hallucinations and identify their spans, defined by character indices, within an answer generated in response to a question. One of the ways to address this is to reformulated the problem as a token classification task, where hallucinated tokens are labeled as 1 and non-hallucinated tokens as 0.

We conducted experiments on transformer-based models to evaluate their performance on this task. Given the limited size of the validation sets (50 labeled data points per language), all base models were trained on our self-generated data (approximately 2k data points per language).

\subsubsection{Data Construction}
It is crucial to have sufficient data to fine-tune the pre-trained language models in order to enhance their performance on specific tasks. However, the training sets provided by the organizers are unlabeled and are thus unable to be used directly for fine-tuning. Therefore, we proposed a semi-automatic approach to construct labeled data ourselves semi-automatically by prompting the state-of-the-art generative language models, specifically GPT-4o.

In order to obtain the data in a manner similar to that of the labeled validation sets provided by the organizers, we explored a combination of few-shot prompting and chain-of-thought (CoT) prompting techniques. Specifically, we utilized several data points in the labeled validations set as learning examples and guided the LLM to infer hallucinated content based on the spans marked by human annotators in those samples. 

We also optimized the pipeline of data construction by prompting the LLM to identify the hallucinated words first and automatically convert the tokens to spans using a simple algorithm afterwards. This additional step significantly increases the quality of generated annotations, as the LLM often misidentified the span boundaries of the hallucinated words. We ultimately constructed 2,371 data points in English and 2000 data points each in the other 12 languages we participated in.


\subsubsection{Base Models}
For fine-tuning, we experimented with a diverse set of pre-trained transformer-based models \cite{vaswani2017attention}, starting with monolingual architectures and later transitioning to multilingual models, which support all 13 target languages, to achieve broader language coverage. Due to time constraints, we focused on English monolingual models for comparison with the performance of multilingual models. Details of the base models used can be found in Table \ref{tab:monolingual_models} for monolingual models and Table \ref{tab:multilingual_models} for multilingual models.

\begin{table*}[h]
    \centering
    \begin{tabular}{|l|l|p{6.5cm}|}
        \hline
        \textbf{Model} & \textbf{Reference} & \textbf{Description} \\ \hline
        \centering\texttt{google-bert/bert-base-cased} & \cite{DBLP:journals/corr/abs-1810-04805} & A bidirectional transformer trained on English text, case-sensitive, widely used for various NLP tasks. \\ \hline
        \texttt{deepset/roberta-base-squad2} & \cite{deepset_roberta_squad2} & A RoBERTa-based model fine-tuned on the SQuAD 2.0 dataset for question answering. \\ \hline
        \texttt{google/flan-t5-base} & \cite{https://doi.org/10.48550/arxiv.2210.11416} & A text-to-text transformer model trained with instruction tuning for better generalization across NLP tasks. \\ \hline
        \texttt{microsoft/deberta-v3-base} & \cite{he2021debertav3} & An improved DeBERTa model with enhanced efficiency and representation learning. \\ \hline
    \end{tabular}
    \caption{Monolingual (English) base models for fine-tuning}
    \label{tab:monolingual_models}
\end{table*}

\begin{table*}[h]
    \centering
    \renewcommand{\arraystretch}{2}  % Increase row height
    \begin{tabular}{|l|l|p{5.8cm}|}
        \hline
        \multicolumn{1}{|p{5.3cm}|}{\parbox[t]{6cm}{\texttt{google-bert/} \\ \texttt{bert-base-multilingual-cased}}} & \cite{DBLP:journals/corr/abs-1810-04805} & A multilingual variant of BERT trained on 104 languages with case sensitivity. \\ \hline
        \multirow{2}{*}
        {\makecell{\texttt{google/umt5-small} \\ \texttt{google/umt5-base}}} & \multirow{2}{*}{\cite{chung2023unimax}} & \multirow{2}{5.8cm}{Multilingual variant of T5 optimized for understanding and generating text across 102 languages.} \\ \cline{1-1}
        & & \\ \hline
        \multirow{2}{*}{\makecell{\texttt{FacebookAI/xlm-roberta-base} \\ \texttt{FacebookAI/xlm-roberta-large}}} & \multirow{2}{*}{\cite{DBLP:journals/corr/abs-1911-02116}} & \multirow{2}{5.8cm}{Multilingual transformer models trained on 94 languages, designed for cross-lingual understanding.} \\ \cline{1-1}
        & & \\ \hline
    \end{tabular}
    \caption{Multilingual base models for fine-tuning}
    \label{tab:multilingual_models}
\end{table*}





\subsection{Performance Evaluation}
The goal of this task is to develop systems that replicate human annotators in detecting hallucinations, identifying spans within a text where hallucinated content occurs. A system's capability to capture hallucination spans is assessed along two main dimensions: (i) the overlap between the system’s predicted hallucination spans and human annotations, and (ii) the alignment in reasoning between the system and human annotators, reflected in the correlation between the confidence of system predictions and the agreement of human annotators on hallucination spans.

These two evaluation dimensions were measured using Intersection over Union (IoU) and Correlation (Cor) scores, respectively. 

The IoU score quantifies the overlap between the predicted hard labels and reference hallucination spans by dividing the size of their intersection by the size of their union. If neither the prediction nor the reference contains hallucinations, the score is set to 1.0.

The Cor score quantifies the agreement between predicted soft labels and reference confidence levels, which are computed as the fraction of annotators who labeled a span as hallucinated. This agreement is measured using Spearman’s rank correlation \cite{spearman1987proof}. The score ranges from -1 to 1, where 1 indicates perfect agreement, 0 signifies no correlation, and -1 represents complete disagreement.



\section{Experimental Setup}
\subsection{Prompt Engineering}
[add content]

\subsection{Model Fine-Tuning}
The fine-tuning procedure began with pre-processing the training data, aligning tokens with binary labels to indicate hallucination. The first stage of fine-tuning was conducted using our auto-generated training data. This step is followed by a second round of fine-tuning using the labeled validation sets provided by the task's organizers, either with all available sets (for multilingual models) or the validation set corresponding to the specific test language (for both monolingual and multilingual models).
Model performance was assessed for both fine-tuning stages. The experimental architecture is illustrated in Figure XX.

[add flow chart showing these steps]

\subsubsection{Data Preprocessing}
\textbf{Label Alignment}  
\label{sec:Label-Alignment}
In all labeled datasets used for fine-tuning, hallucination spans are provided in the format \texttt{[start\_index, end\_index]}. We leveraged this span information to automatically generate labels for each token before feeding the training data into the base models. After tokenization, labels are assigned based on each token’s \texttt{offset\_mapping}: tokens whose start and end indices fall within any hallucination span receive a label of 1, while all others are labeled 0.

\textbf{Data Split}  
We used a 9:1 training-validation split on our self-generated labeled data, resulting in approximately 1,800 training samples and 200 validation samples for each of the 13 languages. These included nine announced target languages: Arabic, German, English, Spanish, Finnish, French, Hindi, Italian, and Chinese, as well as four surprise languages: Czech, Catalan, Basque and Farsi, which were revealed only after the test set was released by the organizers.

We evaluated our models' performance using the labeled validation sets provided by the organizers after fine-tuning. We chose these sets as test data due to (i) their high-quality hallucination span annotated by human annotators and (ii) their likely similarity to the data used by the organizers for final evaluation. In contrast, our automatically generated labeled data were not reviewed by native speakers and may be of lower quality. However, since labeled validation sets were not available for the four surprise languages, we generated 50 labeled data points for each of these languages for testing purpose using the same method as our training data.


\subsubsection{Hyperparameters for Tokenization and Training}
For the fine-tuning procedure, the same set of hyperparameters for both tokenization and training was applied to all base models.  

\textbf{Tokenizer Parameters}  
Only the generated answers (\texttt{model\_output\_text}) were tokenized as input for model fine-tuning, while the inquiries (\texttt{model\_input}) were not used. Tokenization included \texttt{padding} with a maximum length of \texttt{128} tokens and the application of \texttt{truncation}. Additionally, \texttt{return\_offsets\_mapping} was set to \texttt{True}, as the offset mapping is crucial for converting hallucination spans into token-level labels after tokenization. Details on this process are discussed in Section~\ref{sec:Label-Alignment}.


\textbf{Training Parameters}
The training process was configured with a predefined set of parameters. The number of training epochs was set to 10, with a learning rate of 2e-5. Batch sizes were defined as 8 for both training and evaluation (\texttt{per\_device\_train\_batch\_size = 8} and \texttt{per\_device\_eval\_batch\_size = 8}). Additionally, a \texttt{weight\_decay} of 0.01 was applied to the training arguments.






   
\subsection{Benchmarks for Evaluation}
To assess whether fine-tuning improves model performance in hallucination detection, we compare models’ Cor and IoU scores against three benchmarks: (i) the benchmark provided by the organizers, (ii) the base model’s pre-fine-tuning scores (referred to as vanilla scores for convenience), and (iii) the scores from a pre-trained multilingual Named entity recognition (NER) model. 

\textbf{Organizers' Benchmark}
The benchmark provided by the organizers was derived by fine-tuning the multilingual model \texttt{FacebookAI/xlm-roberta-base} \cite{DBLP:journals/corr/abs-1911-02116} using the labeled validation sets they released. These validation sets contained 50 data points per language across 10 languages, excluding the four surprise languages.

Cor and IoU scores were computed for each test language. The base model was trained on validation sets from all languages except the test language, ensuring no test data leakage. This process was repeated for each language, yielding 10 fine-tuned models, each tailored to a specific test language. Fine-tuning was performed using the \texttt{model\_output\_text} and two label types, which classified tokens as either hallucinated or non-hallucinated.




\textbf{Model-Specific Benchmark (Vanilla Scores)}
Evaluating each base model’s scores before fine-tuning helps isolate and quantify the impact of fine-tuning on each model. Since the base models vary in architecture and training data, comparing Cor and IoU scores before and after fine-tuning provides a more accurate measure of its effect. 

The vanilla scores for each base model were obtained using the same methodology and code settings as the fine-tuning process, the only difference being the omission of the training step.

\textbf{NER Benchmark}
A close examination of the sample set shows that a considerable number of hallucinations involve proper nouns, such as names of people, places or organizations. We therefore proposed using a model fine-tuned for NER to make predictions without further fine-tuning, in order to serve as a comparison to the models fine-tuned specifically for this task.

This benchmark was created using the pre-trained NER model \texttt{51la5/roberta-large-NER} \cite{conneau2019unsupervised}, a large multilingual language model supporting all 13 of our target languages. Trained on 2.5TB of filtered CommonCrawl data, this model is an XLM-RoBERTa-large variant fine-tuned on the CoNLL-2003 dataset for English NER. We used the model as-is, without any additional fine-tuning. As a result, it treated all named entities as hallucinations.



\section{Results}


\section{Conclusion}
As our participation in SemEval-2025 Task-3, we proposed three different approaches to predict the spans of hallucination in LLM-generated content, including fine-tuning pretrained language models, named entity recognition and prompting generative large language model. The results show that
Most of our competitive scores across different languages were achieved by fine-tuning the multilingual BERT-based model, XLM-RoBERTa-large, which is proved effective as the scores rank as high as 5th in certain languages.

The fact that the model is exclusively fine-tuned on data constructed by LLM suggests that this approach is a feasible and effective strategy for the task of hallucination span prediction, particularly under circumstances where human-labeled training data is absent. Future directions could include creating and utilizing training data in large quantities, as well as optimizing prompting techniques to obtain higher-quality training data from LLMs. Another viable approach would be to adopt more advanced state-of-the-art models for either data creation or fine-tuning.





\section{Acknowledgments}



\bibliography{custom}
\bibliographystyle{acl_natbib}

\appendix
\section{Appendix}


\end{document}
